\chapter{Introduction\label{cha:chapter1}}

For a plant, photosynthesis controls its growth, carbon uptake and effective biomass formation making its observation across large study areas important for both carbon cycle research and agricultural monitoring. Gross Primary Product (GPP) a measurable proxy for crop yield can be shown as the sum of all leaf photosynthesis \cite{friedlingstein2019global}. Accurately measuring GPP can in turn help find crop yields at parcel, regional levels, but so far it cannot be directly measured by satellites and only remains measurable indirectly through its relationship with net ecosystem carbon exchange through ground based FLUX towers \cite{reichstein2005separation}. Thus only small field level data can be recorded while any regional level statistics cannot be reliably conducted. To alleviate this issue, remote sensing proxies for GPP are needed to expand the spatial scope of monitoring GPP, and one such data is Solar induced fluorescence (SIF). Solar induced fluorescence is an electromagnetic signal emitted by plant chlorophylls during photosynthesis, where a higher SIF signal indicates high sunlight absorption and therefore stronger photosynthesis in stable climatological conditions. Studies like \cite{turner2021extreme, magney2020covariation, frankenberg2011new, joiner2011first} have shown that SIF has a linear relationship with GPP when working with data at satellite scales. Various studies \cite{qiu2022monitoring,song2018satellite,peng2020assessing} have shown SIF to reliably estimate GPP along with vegetation stress, plant heat stress and drought like conditions at large spatial scales. Current SIF recordings have spatial sampling trade offs: OCO-2 and OCO-3 provide comparatively high resolution SIF observations along narrow and widely separated tracks while TROPOMI provides a more spatially complete coverage albeit at a much coarser resolution. Thus it is important to get a high quality SIF product to support the downstream task of GPP and in turn yield monitoring. SIF enhancement methods as shown in Table \ref{tab:sif_enhancement_literature} have enhanced these satellite recorded SIF into higher resolution estimates while ensuring spatio-temporal completeness. Nevertheless, recent enhanced SIF products are limited to approximately 500m or coarser spatial resolution which is still too coarse for isolating crop pure SIF signals in fragmented European crop farms that are usually much smaller than farms in CONUS, Australia, Canada. This thesis therefore tries to address this gap by linking sparse OCO-2 SIF observations over Germany to high resolution Sentinel-2 data and secondary radiation, seasonal, and crop-composition variables represented on a 20~m grid. The resulting 20~m enhanced SIF data is treated as model derived estimates rather than as directly observed or independently validated SIF at 20~m resolution. Finally, the enhanced SIF data is used to test whether crop pure SIF signals can provide more useful information for Bavarian winter wheat yield estimation than raw spatio-temporally sparse OCO-2 footprint data containing mixed vegetation cover. 

\section{Motivation}

Crop specific yield monitoring requires observations that can represent vegetation activity throughout its growing season and also cover large agricultural study regions, remote sensing data from satellites allow this. One such data, Solar induced florescence (SIF) has been shown to be highly effective in monitoring crop yield and GPP  \cite{turner2021extreme, magney2020covariation, frankenberg2011new, joiner2011first}. In the preliminary work of this thesis, raw OCO-2 SIF observations were used to investigate winter wheat yield prediction in Bavaria. Annual crop type rasters were used to filter SIF observations that had high winter wheat coverage and a seasonal SIF dataset was created for Bavaria at NUTS 3, year wise aggregation. However this procedure produced an extremely sparse dataset, with approximately 45-80\% of the required values missing depending on the NUTS-3 region, year, and seasonal period and thus could not be reliably used to estimate winter wheat yields. This preliminary analysis demonstrated that raw OCO-2 data alone is insufficient for spatially complete and crop-specific regional yield analysis. A high resolution SIF data was needed as the average farm size in Bavaria is around 36 hectares (0.36 km$^2$) (calculated from the crop type rasters) and are often highly irregular in shape. Therefore higher resolution corresponds to better isolation of crop specific SIF signal. 

Current raw SIF datasets have trade offs between spatial resolution, temporal coverage and measurement noise. OCO-2 provides relatively high resolution SIF observations in narrow and widely separated satellite scan tracks, while spatio-temporally complete SIF products like TROPOMI, GOME-2 have much coarser resolutions and cannot provide crop pure SIF signals at the field scale. Various enhanced SIF products have been constructed as shown in Table \ref{tab:sif_enhancement_literature}, but the best models are limited to 500m to 1km resolution. At these scales, one SIF pixel can overlap with multiple crops especially for small, irregular shaped farms making it difficult to match them accurately with the fixed grid cells of existing SIF products. A finer SIF representation is therefore needed to isolate winter wheat SIF signals and reduce contamination from neighbouring land / crop cover types. Out of the existing remote sensing data for vegetation indices, optical reflectance bands (prime predictors for SIF enhancement), Sentinel-2 provides the highest resolution at 10m, 20m, and 60m for the 13 reflectance bands. Combining this data with radiation, climatological, seasonal and crop composition statistics can offer a reliable way to estimate SIF at high resolution in locations and time periods without direct OCO-2 retrievals. This highlights the need to develop a more spatially and temporally complete high resolution SIF product and subsequently test whether crop pure enhanced SIF can better estimate crop yields than raw coarse resolution SIF observations. 

\begin{figure}[htbp]
    \centering
    \includegraphics[width=\textwidth]{img/graphical_abstractv1.pdf}
    \caption{Graphical overview of the proposed SIF enhancement and crop yield prediction workflow.}
    \label{fig:graphical_abstract}
\end{figure}

\section{Objective}

The main objective of this thesis is to develop and evaluate a machine learning workflow for enhancing sparse OCO-2 SIF observations by combining OCO-2 supervision with high resolution spectral information as well as radiation, seasonal, and crop composition predictors. Different model architectures like tree based, multi layer perceptron (MLP) and convolutional models are compared to examine how model architecture and input data representation affect SIF estimation accuracy and the spatial distribution of the resulting output enhanced data. The study also tests whether aggregating noisy OCO-2 observations provides a more stable supervision than modelling individual footprints directly as single or multi footprint representations. The data preparation and modelling workflow is prepared in such a way that it can later be applied to a much larger spatial and temporal scope for completing high resolution SIF data. Because complete national multi year processing is computationally intensive, the thesis focuses on producing and evaluating the modelling approach, data-preparation procedures, and reproducible code rather than generating a complete Germany scale product. The proposed modelling approaches are benchmarked against relevant SIF reconstruction \& downscaling methods, with capabilities and limitations highlighted in Section \ref{sec:comparison}. The best model is then selected for extracting winter wheat pure seasonal SIF signals for select Bavarian NUTS 3 regions and agricultural timeframes. Finally, the study identifies whether this crop pure representation adds predictive value relative to raw coarse SIF footprints with incomplete and mixed vegetation coverage.

\section{Thesis Scope}

The main geographic scope for the thesis study is Germany, but instead of a full national coverage select Military Grid Reference System (MGRS) tiles (100km by 100km) covering the main wheat producing states of Bayern, Niedersachsen, Nordrhein-Westfalen, Rheinland-Pfalz, Sachsen-Anhalt, Mecklenburg-Vorpommern, Schleswig-Holstein, Hessen. This spatial scope limitation is needed because multi year Sentinel-2 data processing requires substantial data storage and computational resources. Majority of the SIF enhancement studies in Table \ref{tab:sif_enhancement_literature} use MODIS based reflectance data which can range from 250-500m resolution. In comparison, a 20m resolution raster contains approximately 156 times more pixels than a 250m raster and 625 times more than a 500m raster. Across Germany’s approximately \(357{,}600~\mathrm{km}^2\), one complete raster layer would contain about 1.43 million pixels at 500m, 5.72 million at 250m, 894 million at 20m, and 3.58 billion at 10m and this is only for one band, therefore the data storage and computation can increase drastically. The chosen 13 MGRS tiles of Sentinel-2 data is around 500GB before further model ready processing. The temporal scope is 2019-2024 because the Sentinel-2 L3A WASP monthly archive used in the thesis begins in 2019 and the required annual crop type rasters exist up to 2024. Although several secondary vegetation and radiation datasets cover longer periods, they are restricted to this shared interval to maintain a temporally consistent set of predictor channels. The thesis's SIF enhancement experiments use available SIF observations and predictors from February to July, which is the main growing period for winter wheat in Germany. The yield estimation experiment worked with a more narrow geographic scope and covers Bavarian NUTS 3 regions within five selected MGRS tiles using monthly SIF from March to July, years 2019-2022 for model training, and 2024 as an independent temporal test year. The year 2023 and month February were excluded from the yield experiment because matching optical predictor coverage was incomplete and a complete mean aggregated SIF February season variable had a lot of missing values.

\section{Purpose and goal\label{sec:purpose_and_goal}}

The main purpose of this thesis is to bridge the scale gap between sparse OCO-2 SIF observations and the field scale spatial information required for crop specific agricultural analysis. It aims to determine whether high resolution Sentinel-2 and secondary predictors can distribute footprint or grid supported SIF into accurate spatial patterns that are useful at the scale of fragmented German farmland. The secondary purpose is to determine if restricting enhanced SIF estimates to winter wheat pixels reduces the mixed crop cover problem encountered when raw OCO-2 SIF observations are aggregated for regional yield prediction. The study also aims to clarify how predictor information, spatial context, OCO-2 aggregation, and target preparation influence SIF enhancement performance instead of attributing all differences to model architecture alone. Beyond crop yield estimation, spatially detailed and temporally complete SIF could support the monitoring of plant phenological development, and physiological responses to drought or heat stress conditions \cite{song2018satellite}. But this would require changing the temporal representation of SIF. Rather than representing it as daily or multi-day averages, it would require sub daily SIF observations such as separate morning and afternoon estimates. Eventhough studies like \cite{li2019global, ma2020generation, yu2019high} provide complete enhanced SIF products at \(0.05^\circ\) scale for large timeframes, this would not be computationally feasible for this study given high increase in data information moving from 0.05 to 20m resolution. Instead the study contribution is a reproducible modelling and data preparation workflow that can support later generation of longer and more geographically
complete SIF records. The intended outcome is therefore an evaluated high resolution SIF enhancement worflow, a transparent account of its validation limits and evidence on whether crop pure enhanced SIF provides a useful direction for future yield research.

\section{Outline\label{sec:outline}}

The remainder of this thesis is organized into six chapters that progress from theoretical foundations and related research for the study, research questions, experimental workflow design, experiments conducted, their findings and conclusions.
\\
\\
\textbf{Chapter \ref{cha:chapter2}} introduces the theoretical background required to understand photosynthesis and light use theory, SIF to GPP relationship, various remote sensing predictors. It also explains the modelling concepts, spatial supports and evaluation metrics.
\\
\\
\textbf{Chapter \ref{cha:chapter3}} reviews existing approaches for SIF reconstruction, gap filling and downscaling methods followed by remote sensing methods for crop yield and GPP estimation. The chapter also compares main advantages and limitations for a select set of studies and identifies the research gap addressed by this thesis. 
\\
\\
\textbf{Chapter \ref{cha:chapter4}} defines the direct usability problem created by the sparse spatial sampling of OCO-2 and the difficulty of obtaining crop specific SIF information. It then presents the research questions concerning SIF prediction, model representation, spatial aggregation, generalization, winter-wheat yield estimation, effect of sample density on field scale reconstructed SIF uncertainty and transferability of enhancement framework to other c3 crops. 
\\
\\
\textbf{Chapter \ref{cha:chapter5}} describes the research design, study area, datasets, preprocessing procedures, and construction of the model-ready representations. It also documents the evaluated model architectures, training and validation strategies, experimental environment, reproducibility measures, and how the evaluation metrics are applied. 
\\
\\
\textbf{Chapter \ref{cha:chapter6}} presents and discusses the experimental results for SIF prediction, model architecture, input representation, spatial aggregation, and model generalization to previously unseen study areas. It also evaluates spatial and seasonal performance differences, compares crop pure and raw SIF for winter wheat yield prediction, and relates the findings to existing enhancement techniques.
\\
\\
\textbf{Chapter \ref{cha:chapter7}} summarizes the experimental workflow and the principal findings and considers the extent to which the research questions were answered. It also discusses dissemination, practical problems and limitations encountered during the study, and possible directions for extending the SIF enhancement and crop monitoring workflow.